\documentclass[a4paper, 11pt]{article}

\usepackage[swedish]{babel}
\usepackage[utf8]{inputenc}
\usepackage[T1]{fontenc}
\usepackage{lmodern}

\usepackage{tikz}

\tikzset{
  excl/.style = {circle, draw=black, inner sep=0pt, minimum size=4pt}
}

\tikzset{
  incl/.style = {circle, draw=black, fill=black, inner sep=0pt, minimum size=4pt}
}


\usepackage{pgfplots}


\usepackage{graphicx}
\usepackage{caption}
\usepackage{subcaption}

\usepackage{enumitem}

\usepackage{multicol}

\usepackage{amsmath}


\begin{document}

\title{
    \textbf{MM2001: Inläming 0}
}
\author{Johan Montelius}
\date{HT 2026}

\maketitle



\section*{i)  Persson-Böiers 1.6.1-2}


\subsection*{a) Förenkla uttrycket:}


$$ \frac{\sqrt[3]{ a^6}}{(\sqrt[3]{\sqrt[4]{a^5}})^6} \times \sqrt{a}$$.


\subsection*{b) Visa att $x = 9/4$ är en lösning till ekvationen: 

$$x^{x \sqrt{x}} \eq (x\sqrt{x})^x)$$.

\subsection{ c) Har du nu löst ekvation ovan?}

Nej - jag har visat att den gäller för $x = 9/4$ men inte löst ut ett generellet uttryck som ger värdet för $x$.


\section*{ii) Persson-Böiers 1.6.1}

\subsection*{a)} Uttryck formelrad (27) med ord.


\subsection*{a Vad händer i (27) om $\alpha \leq 0$?}

  

\end{document}


